% ===== AI-DRAFTED (pending author review) =====
% Auto-generated from survey/content.html by scripts/sync_html_to_tex.js — do not edit by hand.
\section{Preliminaries and Notation}\label{sec:prelim}
\aibanner

We collect the notation used throughout. An instance consists of a set of $n$ data points $P$ in a metric space $(X,d)$, where $d$ obeys the triangle inequality. A \emph{clustering} partitions $P$ into $k$ groups $\mathcal{C}=\{C_1,\dots,C_k\}$; in \emph{centroid-based} formulations each cluster $C_i$ is represented by a center $c_i$ from a candidate set $F\subseteq X$ (often $F=P$), and each point is assigned to its nearest open center.

\paragraph{Clustering objectives.} The three canonical objectives are unified by the $\ell_p$ cost of a center set $C$, $\mathrm{cost}_p(P,C)=\big(\sum_{x\in P} d(x,C)^p\big)^{1/p}$ with $d(x,C)=\min_{c\in C} d(x,c)$. Here $p=1$ is \textbf{$k$-median}, $p=2$ is \textbf{$k$-means} (sum of squared distances), and $p\to\infty$ is \textbf{$k$-center}, minimizing $\max_{x} d(x,C)$. All three are NP-hard, so the literature targets \emph{approximation} guarantees.

\paragraph{Protected groups.} Group fairness is defined relative to $\ell$ protected groups $P_1,\dots,P_\ell\subseteq P$ from a sensitive attribute. Groups may be \emph{disjoint} or \emph{overlapping}; $\Delta$ denotes the maximum number of groups any point belongs to ($\Delta=1$ when disjoint).

\paragraph{Approximation.} An $\alpha$-approximation returns a feasible (fair) solution of cost at most $\alpha$ times optimal. Many results are \emph{bicriteria}: an $(\alpha,\beta)$-approximation achieves cost within $\alpha$ of optimal while relaxing fairness by a slack $\beta$. Noting \emph{which} quantity is relaxed --- cost or fairness --- is essential when comparing results.

The taxonomy in Figure~1 organizes the notions surveyed below by \emph{where} fairness is enforced: within clusters, among centers, across group costs, or at the level of individuals and coalitions.
